\documentclass[reqno, 12pt]{amsart}

\usepackage{amssymb}
\usepackage{amsthm}
\usepackage{amsmath}
\usepackage{amsxtra}
\usepackage{mathrsfs}
\usepackage{comment}
\usepackage{graphicx}
\usepackage{placeins}
\usepackage{multicol}
\usepackage{bbm}
\usepackage{stmaryrd}
%\usepackage{enumerate}
\usepackage[inline, shortlabels]{enumitem}
\usepackage{mathdots}
\usepackage{colonequals}
\usepackage{hyperref}


\usepackage{calc}
\usepackage{tikz}
\usetikzlibrary{arrows,calc,automata,shadows,backgrounds,positioning,intersections,fadings,decorations.pathreplacing,shapes,matrix}
\usepackage{tikz-cd}

\usepackage{fullpage}

\newtheorem{thm}{Theorem}[section]
\newtheorem{lem}[thm]{Lemma}
\newtheorem{cor}[thm]{Corollary}
\newtheorem{conj}[thm]{Conjecture}
\newtheorem{prop}[thm]{Proposition}
\theoremstyle{definition}  
\newtheorem{defn} [thm] {Definition} 
\newtheorem{claim}[thm]{Claim}
\newtheorem{example} [thm] {Example}
\newtheorem{rem} [thm] {Remark}

\newenvironment{problab}[1]
{\noindent\textbf{Problem #1}.}
{\vskip 6pt}

\newenvironment{exolab}[1]
{\noindent\textbf{Exercise #1}.}
{\vskip 6pt}

\theoremstyle{remark}
\newtheorem*{soln}{Solution}

\newcommand{\C}{\mathbb C}
\renewcommand{\H}{\mathbb H}
\newcommand{\D}{\mathbb D}
\newcommand{\F}{\mathbb F}
\newcommand{\Q}{\mathbb Q}
\newcommand{\R}{\mathbb R}
\newcommand{\Z}{\mathbb Z}
\newcommand{\T}{\mathbb T}
\renewcommand{\P}{\mathcal P}
\renewcommand{\O}{\mathcal O}
\newcommand{\m}{\mathfrak m}
\newcommand{\A}{\mathbb A}
\renewcommand{\SS}{\mathbb S}
\renewcommand{\L}{\mathcal L}

\newcommand{\B}{\mathcal B}
\newcommand{\E}{\mathcal E}

\newcommand{\SSpan}{\text{span}}
\newcommand{\la}{\langle}
\newcommand{\ra}{\rangle}

\DeclareMathOperator{\lcm}{lcm}
\DeclareMathOperator{\img}{img}
\DeclareMathOperator{\Ann}{Ann}
\DeclareMathOperator{\Tor}{Tor}
\DeclareMathOperator{\tr}{tr}
\DeclareMathOperator{\Hom}{Hom}
\DeclareMathOperator{\End}{End}
\DeclareMathOperator{\Aut}{Aut}
\DeclareMathOperator{\Gal}{Gal}
\DeclareMathOperator{\sgn}{sgn}
\DeclareMathOperator{\ord}{ord}
\DeclareMathOperator{\ad}{ad}
\DeclareMathOperator{\coker}{coker}
\DeclareMathOperator{\rank}{rank}
\DeclareMathOperator{\minpoly}{minpoly}

\DeclareMathOperator{\id}{id}

\usepackage{mathpazo}

\everymath{\displaystyle}

\tikzset{commutative diagrams/.cd, arrow style = tikz, diagrams = {>=latex}}

\newenvironment{amatrix}[1]{%
  \left(\begin{array}{@{}*{#1}{c}|c@{}}
}{%
  \end{array}\right)
}

\usepackage{abstract}
\renewcommand{\abstractname}{}    % clear the title
\renewcommand{\absnamepos}{empty} % originally center
%\setlength{\abstitleskip}

\begin{document}
%\renewcommand{\abstractname}{\vspace{-\baselineskip}}

\title{18.700 Problem Set 7}
%\author{Évariste Galois}
\dedicatory{Due Wednesday, November 6 at 11:59 pm on \href{https://canvas.mit.edu/courses/27315}{Canvas}}
%\thanks{}
%\contrib{}
%\address{}

\maketitle

\begin{abstract}
  \noindent Collaborated with:\\
  Sources used: 
\end{abstract}
\vskip 0.5cm

Let $V$ be an inner product space.

\begin{problab}{1} (6A \#9, 11) (10 points)
  \begin{enumerate}[(a)]
    \item 
      Suppose $u,v \in V$, and $\|u\| = 1 = \|v\|$ and $\la u, v \ra = 1$. Prove that $u = v$.
    \item
      Find vectors $u,v \in \R^2$ such that $u$ is a scalar multiple of $(1,3)$, $v$ is orthogonal to $(1,3)$, and $(1,2) = u+v$.
  \end{enumerate}
\end{problab}

%\begin{soln}
%\end{soln}


\begin{problab}{2} (6A \#16) (5 points)
  Given $x,y \in V$, we defined the angle between $x$ and $y$ to be
  \begin{equation*}
    \angle(x,y) \colonequals \arccos\left( \frac{\la x, y \ra}{\|x\| \|y\|} \right) \, .
  \end{equation*}
  Explain why this definition makes sense using the Cauchy-Schwarz inequality.
\end{problab}

\begin{problab}{3} (6B \#1) (6 points)
  Suppose $e_1, \ldots, e_m$ is a list of vectors in $V$ such that
  \begin{equation*}
    \|a_1 e_1 + \cdots + a_m e_m\|^2 = |a_1|^2 + \cdots + |a_m|^2
  \end{equation*}
  for all $a_1, \ldots, a_m \in \F$. Show that $e_1, \ldots, e_m$ is an orthonormal list. (This exercise provides a converse to 6.24 in the textbook.)
\end{problab}

\begin{problab}{4} (6B \#7) (6 points)
  Suppose $T \in \L(\R^3)$ has an upper-triangular matrix with respect to the basis $(1,0,0), (1,1,1), (1,1,2)$. Find an orthonormal basis of $\R^3$ with respect to which $T$ has an upper-triangular matrix.
\end{problab}

\begin{problab}{5} (6C \#3) (8 points)
  Suppose $U$ is the subspace of $\R^4$ defined by
  \begin{equation*}
    U \colonequals \SSpan\left( (1,2,3,-4), (-5,4,3,2) \right) \, .
  \end{equation*}
  Find an orthonormal basis of $U$ and an orthonormal basis of $U^\perp$.
\end{problab}

\begin{problab}{6} (6C \#4) (10 points)
  Suppose $e_1, \ldots, e_n$ is a list of vectors in $V$ with $\|e_k\| = 1$ for each $k = 1, \ldots, n$, and
  \begin{equation*}
    \|v\|^2 = |\la v, e_1 \ra|^2 + \cdots + |\la v, e_n \ra|^2
  \end{equation*}
  for all $v \in V$. Prove that $e_1, \ldots, e_n$ is an orthonormal basis of $V$. (This provides a converse to 6.30(b) in the textbook.)
\end{problab}

\end{document} 	
